\section{From Hodge cancellation to physical observables}
\label{sec:r7-hodge-bridge}

Let $u$ be canonical Hodge coordinates for an isolated spectral branch
$E(u)$ and let $\|u\|_G$ be the Hodge norm.

\begin{theorem}[Content and limits of Hodge cancellation]
\label{thm:r7-hodge-bridge}
Assume $dE(0)=0$, $\nabla_G^2E(0)=0$, and
$\|D^3E(u)\|_{G\to\mathbb R}\le M_3$ on $\|u\|_G\le r$.  Then
\begin{equation}
 |E(u)-E(0)|\le M_3\|u\|_G^3/6.
 \label{eq:r7-hodge-cubic}
\end{equation}
If inversion symmetry also removes the cubic tensor and
$\|D^4E\|\le M_4$, the right side improves to
$M_4\|u\|_G^4/24$.

Suppose, in addition, that the Riesz-compressed operator satisfies the
stronger norm statement
\begin{equation}
 \|(H_{\rm eff}(u)-E(0))P(u)\|\le\epsilon(u).
 \label{eq:r7-effective-flatness}
\end{equation}
For every normalized compact wavepacket $\psi\in\operatorname{Ran}P(u)$,
\begin{align}
 \operatorname{Var}_\psi(H)&\le\epsilon(u)^2,\\
 \|e^{-itH}\psi-e^{-itE(0)}\psi\|&\le |t|\epsilon(u),\\
 \langle W_\ell\rangle_t&\le t^2\epsilon(u)^2
 \quad\text{when }W_\ell\psi=0,\\
 -\eta\Im\langle\psi,(E(0)+i\eta-H)^{-1}\psi\rangle
 &\ge\eta^2/(\eta^2+\epsilon(u)^2).
 \label{eq:r7-hodge-observables}
\end{align}
Thus complete Hodge cancellation controls local spectral and dynamical
diagnostics only after an isolated-branch, projector-overlap, and operator-norm
flatness estimate is supplied.
\end{theorem}

\begin{proof}
Taylor's theorem proves \eqref{eq:r7-hodge-cubic}.  Under
\eqref{eq:r7-effective-flatness}, the spectral measure of $\psi$ lies in the
$\epsilon(u)$ interval about $E(0)$, which proves the variance and Poisson
kernel estimates.  Duhamel's formula proves the propagator estimate.  If
$W_\ell$ is a positive contraction annihilating the reference state, its
expectation is bounded by the squared norm of the orthogonal departure,
giving the third line.
\end{proof}

Hodge cancellation alone does \emph{not} determine the local return or the
real-space spreading coefficient.  For example,
\begin{equation}
 \left.\frac{d^2}{dt^2}\langle X^2\rangle_t\right|_{t=0}
 =-\langle\psi,[H,[H,X^2]]\psi\rangle
\end{equation}
depends on real-space commutators, whereas $\nabla_G^2E$ differentiates a
representation-space branch.  Models can share a zero Hodge Hessian and have
different $[H,X]$.  In the frozen anchor the missing implication is supplied
not by rhetoric but by the exact first-shell identity $H_{+,1}=E_*I$ and the
full-kernel norm bound $\varepsilon_A$.  Substitution
$\epsilon=\varepsilon_A$ yields Proposition~\ref{prop:r7-anchor-margin}.

For a local state with Riesz weight
$\alpha=\langle\psi,P\psi\rangle<1$, the complement contributes explicitly;
a safe variance bound is
\begin{equation}
 \operatorname{Var}_\psi(H)
 \le 2\epsilon^2+8M^2(1-\alpha).
\end{equation}
This records the active-overlap hypothesis rather than silently replacing it
by one.
